\medskip






L'espace est rapporté à un repère orthonormé \Oijk.

\medskip

\begin{list}{\textbullet}{On considère:}
\item les points A$(-1~;~-2~;~3)$, B$(1~;~-2~;~7)$ et C$(1~;~0~;~2)$;
\item la droite $\Delta$ de représentation paramétrique : \quad $\left\{\begin{array}{l}x=1-t \\ y=2 \\ z=-4+3 t\end{array}\right.$, où $t \in \R$;
\item le plan $\mathcal{P}$ d'équation cartésienne :\quad $3 x+2 y+z-4=0$;
\item le plan $\mathcal{Q}$ d'équation cartésienne : \quad $-6 x-4 y-2 z+7=0$.
\end{list}

\begin{enumerate}[itemsep=5mm]
	\item Lequel des points suivants appartient au plan $\mathcal{P}$ ?

	\begin{tabularx}{\linewidth}{*{4}{>{\bfseries}p{5mm}@{}X}}
		a.& $R(1~;~-3~;~1)$;&
		b.& $S(1~;~2~;~-1)$;&
		c.& $T(1~;~0~;~1)$;&
		d.& $U(2~;~-1~;~1)$.\\
	\end{tabularx}

\medskip

\begin{tabular}{@{\hspace*{0.05\linewidth}} || p{0.93\linewidth}}
$3x_{\text T} +2y_{\text T} +z_{\text T}-4
= 3+0+1-4=0$ donc $\text T \in \mathcal{P}$
\hfill\textbf{Réponse c.}
\end{tabular}



	\item Le triangle ABC est :

	\begin{tabularx}{\linewidth}{*{2}{>{\bfseries}p{5mm}@{}X}}
		a.& équilatéral ;&
		b.& rectangle isocèle ;\\
		c.& isocèle non rectangle ;&
		d.& rectangle non isocèle.\\
	\end{tabularx}

\medskip

\begin{tabular}{@{\hspace*{0.05\linewidth}} || p{0.93\linewidth}}
Les vecteurs $\vectt{AB}$ et $\vectt{AC}$ ont respectivement pour coordonnées
$\begin{pmatrix}
2\\0\\4
\end{pmatrix}$
et
$\begin{pmatrix}
2\\2\\-1
\end{pmatrix}$.

Donc $\vectt{AB}\cdot\vectt{AC} = 2\times 2+0\times 2+4\times (-1)=0$ donc $\vectt{AB}\perp \vectt{AC}$ donc ABC est rectangle en A.

$\text{AB}^2=2^2+0^2+4^2=20$ et $\text{AC}^2=2^2+2^2+(-1)^2=9$, donc $\text{AB}\neq \text{AC}$

Le triangle ABC n'est pas isocèle.
\hfill\textbf{Réponse d.}
\end{tabular}

\medskip

	\item La droite $\Delta$ est :

	\begin{tabularx}{\linewidth}{*{2}{>{\bfseries}p{5mm}@{}X}}
		a.& orthogonale au plan $\mathcal{P}$ ;&
		b.& sécante au plan $\mathcal{P}$ ;\\
		c.& incluse dans le plan $\mathcal{P}$ ;&
		d.& strictement parallèle au plan $\mathcal{P}$.\\
	\end{tabularx}

\medskip

\begin{tabular}{@{\hspace*{0.05\linewidth}} || p{0.93\linewidth}}
La droite $\Delta$ a pour vecteur directeur
$\vect v
\begin{pmatrix}
-1\\0\\3
\end{pmatrix}$,
et le plan $\mathcal{P}$ a pour vecteur normal
$\vect n
\begin{pmatrix}
3\\2\\1
\end{pmatrix}$.

$\vect v\cdot \vect n = -1\times 3 +0\times 2+3\times 1 = 0$ donc $\vect v\perp \vect n$, et donc $\Delta$ est parallèle à $\mathcal{P}$.

Le point H de coordonnées $(1\;;\;2\;;\;-4)$ appartient à $\Delta$.

$3x_{\text H} +2y_{\text H} +z_{\text H}-4 
=3\times 1 + 2\times 2 -4-4=-1\neq 0$ donc $\text H \notin\mathcal{P}$.

La droite $\Delta$ n'est pas incluse dans le plan $\mathcal{P}$.
\hfill\textbf{Réponse d.}
\end{tabular}

\medskip

	\item On donne le produit scalaire $\vect{\mathrm{BA}} \cdot \vect{\mathrm{BC}}=20$.

	Une mesure au degré près de l'angle $\widehat{\mathrm{ABC}}$ est :

	\begin{tabularx}{\linewidth}{*{4}{>{\bfseries}p{5mm}@{}X}}
		a.& $34^{\circ}$;&
		b.& $120^{\circ}$;&
		c.& $90^{\circ}$;&
		d.& $0^{\circ}$.\\
	\end{tabularx}

\medskip

\begin{tabular}{@{\hspace*{0.05\linewidth}} || p{0.93\linewidth}}
Le triangle ABC est rectangle en A donc l'angle $\widehat{\mathrm{ABC}}$ est un angle aigu. 
\hfill\textbf{Réponse a.}
\end{tabular}

\medskip

	\item L'intersection des plans $\mathcal{P}$ et $\mathcal{Q}$ est :

	\begin{tabularx}{\linewidth}{*{2}{>{\bfseries}p{5mm}@{}X}}
		a.& un plan ;&
		b.& l'ensemble vide ;   \\
		c.& une droite ;&
		d.& réduite à un point.\\
	\end{tabularx}
	
\medskip

\begin{tabular}{@{\hspace*{0.05\linewidth}} || p{0.93\linewidth}}
Le plan $\mathcal{P}$ a pour vecteur normal
$\vect n
\begin{pmatrix}
3\\2\\1
\end{pmatrix}$
et le plan $\mathcal{Q}$ a pour vecteur normal
$\vect{n'}
\begin{pmatrix}
-6\\-4\\-2
\end{pmatrix}$.
$\vect{n'}=-2\vect n$ donc les plans sont parallèles.

Le point T appartient à $\mathcal{P}$.

$-6x_{\text T} -4y_{\text T} -2z_{\text T} +7
= -6\times 1 -4\times 0 -2\times 1+7
=-1\neq 0$ donc $\text T \notin \mathcal{Q}$

Les deux plans sont donc strictement parallèles.
\hfill\textbf{Réponse b.}
\end{tabular}

\end{enumerate}
